\uuid{Rp7E}
\titre{Étude de la classe $\mathcal{C}^1$ de $f(x,y)=e^{-1/(x^2+y^2)}$}
\niveau{L2}
\module{Analyse}
\chapitre{Fonctions de plusieurs variables}
\sousChapitre{Continuité et différentiabilité}
\theme{Dérivées partielles, continuité, croissances comparées, classe $\mathcal{C}^1$}
\auteur{Maxime Nguyen}
\datecreate{2026-03-24}
\organisation{amscc}
\difficulte{4}

\contenu{
\texte{
  On considère la fonction $f$ définie sur $\mathbb{R}^2$ par
  \[
    f(x,y) = e^{-\frac{1}{x^2+y^2}} \;\text{ si } (x,y) \neq (0,0), \qquad f(0,0) = 0.
  \]
  Cette fonction est-elle de classe $\mathcal{C}^1$ ?
}
\begin{enumerate}
  \item \question{Calculer les dérivées partielles de $f$ pour $(x,y) \neq (0,0)$.}
    \reponse{
      \[
        \frac{\partial f}{\partial x}(x,y) = \frac{2x}{(x^2+y^2)^2}\,e^{-\frac{1}{x^2+y^2}}, \qquad \frac{\partial f}{\partial y}(x,y) = \frac{2y}{(x^2+y^2)^2}\,e^{-\frac{1}{x^2+y^2}}.
      \]
    }

  \item \question{Calculer les dérivées partielles de $f$ en $(0,0)$.}
    \reponse{
      Pour $h \neq 0$, $f(h,0) = e^{-1/h^2}$. Par croissances comparées, $\displaystyle\lim_{h\to 0}\dfrac{e^{-1/h^2}}{h} = 0$. Donc $\dfrac{\partial f}{\partial x}(0,0) = 0$. De même $\dfrac{\partial f}{\partial y}(0,0) = 0$.
    }

  \item \question{Étudier la continuité des dérivées partielles en $(0,0)$.}
    \reponse{
      En coordonnées polaires :
      \[
        \left|\frac{\partial f}{\partial x}(x,y)\right| = \frac{2|\cos\theta|}{r^3}\,e^{-1/r^2} \leq \frac{2}{r^3}\,e^{-1/r^2}.
      \]
      Par croissances comparées, $\dfrac{1}{r^3}e^{-1/r^2} \to 0$ quand $r \to 0$. Donc $\dfrac{\partial f}{\partial x}$ est continue en $(0,0)$, et de même pour $\dfrac{\partial f}{\partial y}$. Ainsi $f$ est de classe $\mathcal{C}^1$ sur $\mathbb{R}^2$.
    }
\end{enumerate}
}
